\section{Robustness and Research Implications}
\subsection{Observation and Selection Sensitivities}
We exclude 60 positive observations whose stored repository key disagrees with a lexical parse of the source URL, reducing the endpoint universe by 24 repositories (7,171 to 7,147); greedy at 100 covers 5,524, with $P=799$. We do not transfer workflow flags without reacquisition. This check addresses a specific mapping defect and cannot certify remaining remote associations.

Removing all bundled-context nodes leaves every workflow-repository source set unchanged, including all 462,170 identity--repository pairs. Consequently, primary repository $C$, $H$, and $P$ are identical for any fixed selection. Workflow-package incidence changes by 28 pairs for one endpoint, which repository deduplication absorbs; contextual nodes should therefore not be assumed irrelevant to every endpoint definition.

Reversing tie order yields $C=5{,}537$ and $P=793$. Excluding the 1, 5, or 10 largest portfolios gives coverage of 5,522, 5,149, and 5,076. Grouping equal stored email hashes as a hypothetical control unit yields 12,147 groups with greedy coverage at 5,538. These tests reflect eligibility sensitivity rather than verified credential control.

\subsection{From Structural Reach to Outcome Research}
Our findings suggest distinct uses for the quantities. A team seeking broad review coverage can use $C$; the stipulated fixed-prior account-hardening decision uses $H$. Complete modeled source exclusion requires $P$ and may demand complementary selections. Actual budgets also depend on control cost, feasibility, account-specific probabilities, publishing permissions, and adversary adaptation. Our equal-cost fixed-prior experiment does not observe these. Workflow association alone cannot establish credential-harvesting return or a control's causal effect.

We completed prospective acquisition groundwork by sampling 100 identifiers uniformly without replacement from our 25,228-name frame (seed 20261020, frozen before requests). Metadata were captured for 90; eight HTTP 405 responses involve bundled identifiers, while two return HTTP 404. The feature cutoff is 20 October 2026, 00:48:38 Asia/Dhaka, with a 90-day horizon ending 18 January 2027. All 100 outcomes remain pending. This pilot demonstrates feasibility; its size does not establish event yield or predictive power.

Future compromise evaluation must fix surveillance sources, reporting delays, censoring, and family-grouped splits before examining outcomes. Unreported packages remain unlabeled rather than certified benign. Contemporary metadata cannot be retroactively assigned to earlier attacks. We report no trained forecast, accuracy, or empirical compromise probability from this pilot.
